\section{Protocol for comparative evaluation}

A performance evaluation must separate three effects that are easily conflated:
capacity, runtime overhead, and latency hiding. Shadow Trainer may enable a job
under a disk ceiling even when it is slower than framework-local loading. A
prefetch arm may reduce exposed I/O time while leaving total runtime unchanged
because training dominates. Conversely, a warm source or filesystem cache may
appear to accelerate prefetch even when overlap provides no benefit. We define
the following protocol to prevent these cases from becoming one aggregate
``speedup'' number.

\subsection{Arms and controlled variables}

The framework-local arm reads a fully local dataset without Shadow Trainer and
establishes orchestration overhead. The synchronous arm uses atomic staging and
the bounded cache. The prefetch arm changes only whether the next case window is
materialized concurrently. Paired sync and prefetch processes use the same
manifest bytes, case order, seed, initial checkpoint, model options, training
steps, cache capacity, disk floor, driver, CUDA, cuDNN, and PyTorch version.
Each arm starts in a new process to avoid allocator and library state leaking
between measurements.

Cold and warm conditions are separate populations. A cold run starts from an
independently prepared cache state and records every source byte. A warm run is
allowed to reuse the declared case cache but must identify its initial contents.
Operating-system page-cache state is difficult to control safely on a shared
workstation; it is therefore reported as a limitation rather than manipulated
with privileged cache-dropping commands.

\subsection{Measurements}

Primary capacity outcomes are admission, completion, peak allocated and reserved
VRAM, process RSS, swap use, maximum cache occupancy, combined artifact bytes,
and minimum observed disk headroom. Primary recovery outcomes are lost or
repeated step identifiers and equality of final model, optimizer, RNG, and data
position. Performance outcomes are end-to-end time, staging time, bytes moved,
cache hits and misses, exposed I/O stall, steps per second, and GPU utilization.

Every reported performance cell requires at least three independent processes;
five are preferred for short workloads. We report individual observations,
median, interquartile range, and a bootstrap confidence interval for the paired
difference. A mean alone is not used for small, skewed I/O samples. Hardware
telemetry sampling must be identical across arms and sufficiently light that its
overhead is either measured or held constant.

\subsection{Exactness gate}

Latency hiding is eligible for a performance comparison only after the complete
observable result passes an exactness audit. The audit checks ordered case and
step identifiers, finite losses, final model and optimizer tensors, RNG states,
checkpoint position, manifest digest, and boundary hashes. If deterministic
algorithms cannot make a particular kernel bitwise stable, the protocol must
predeclare numerical tolerances and justify them; a tolerance cannot be selected
after observing the difference. The current product takes the stricter position:
prefetch remains experimental because the full pair has not completed.

\subsection{Capacity experiment}

The principal paper result should use a remote dataset whose declared size
exceeds the cache and the disk allocation available to the job. Capacity is
demonstrated only if all requested cases are processed, maximum cache and
artifact occupancy remain within contract, every promoted file passes integrity
validation, and the final disk floor is respected. Comparing total dataset size
with physical disk size is unnecessary and potentially misleading; the relevant
quantity is the storage allocation made available under the experimental
contract.

The same experiment should be repeated for NIfTI 3D, ResNet-18, ResNet-50, and
a real adjacent-slice 2.5D pipeline. The current proxy ResNet runs validate the
adapter boundary and topology memory footprint but do not replace real decoding,
augmentation, or data-loader behavior. Together, the capacity and overhead arms
would determine whether the unified contract provides a useful capability and
what performance price it imposes.
